\chapter{Exponential and Logarithm}

\section{Exponential}

\begin{definition}[Exponential representations]\label{def:exp-reps}
\lean{ComputableAnalysis.expPowerSeries}
\lean{ComputableAnalysis.expEuler}
\lean{ComputableAnalysis.ePowerSeries}
\lean{ComputableAnalysis.eEuler}
\lean{ComputableAnalysis.ExpProofs.EPowerSeriesEqEuler}
\uses{def:raw-equiv}
The project compares several algorithms for the exponential function:
\[
  \exp_{\mathrm{series}}(x)
  =
  \sum_{n=0}^{\infty}\frac{x^n}{n!},
\]
\[
  \exp_{\mathrm{Euler}}(x)
  =
  \lim_{m\to\infty}\left(1+\frac{x}{m}\right)^m,
\]
and the inverse of the logarithmic integral once that branch has been
constructed.  Each representation is a rational interval algorithm, and
equality means equivalence of raw real numbers.
\end{definition}

\begin{definition}[Power-series exponential as a raw algorithm]
\label{def:exp-power-series-raw}
\lean{ComputableAnalysis.expPowerSeriesPartialAndTailBound}
\lean{ComputableAnalysis.expPowerSeries}
\uses{def:exp-reps,def:raw-real}
At stage \(n\), the power-series algorithm computes a rational partial sum and
an explicit rational tail bound.  The tail is controlled by a geometric
majorant for the omitted terms, so the returned interval is
\[
  [S_n-R_n,S_n+R_n].
\]
\end{definition}

\begin{definition}[Euler-limit exponential]
\label{def:exp-euler-raw}
\lean{ComputableAnalysis.expEuler}
\uses{def:exp-reps,def:raw-real}
The Euler representation computes the finite rational power
\[
  \left(1+\frac{x}{m}\right)^m
\]
with an explicit display interval at each stage.  The equivalence with the
power-series representation is one of the central estimates for this chapter.
\end{definition}

\section{Logarithm}

\begin{definition}[Logarithm by the integral of \(1/x\)]
\label{def:log-integral}
\lean{ComputableAnalysis.RatFun.oneOverXOnPositiveInterval}
\lean{ComputableAnalysis.exp.LogIntegralInverseBranch}
\uses{def:definite-integral-identity}
On a positive rational interval, the logarithm branch is represented by
\[
  \log x=\int_1^x \frac{\dd t}{t}.
\]
The interval chapter supplies the definite-integral identity package and the
positive-interval domain certificate for \(1/x\).  The exponential function
obtained by inverting this logarithmic integral is then compared with the
power-series and Euler-limit exponentials.
\end{definition}

\begin{definition}[Agreement of exponential representations]
\label{def:exp-representations-agree}
\lean{ComputableAnalysis.exp.representationsAgree}
\lean{ComputableAnalysis.exp.representationsAgree_of_estimates}
\uses{def:exp-reps,def:log-integral}
The intended theorem package says that the power-series exponential, the
Euler-limit exponential, and the inverse-log-integral exponential agree on
their common domains.  The proof can be organized by direct estimates or by
showing that the competing representations solve the same effective
differential equation with the same initial value.
\end{definition}
